\section{Counting and visiting collision witnesses}\label{sec:witnesses}

We now bound the probability of the marked path unions furnished by
Lemma~\ref{lem:witness}. The required precision is
$o(D^{-1})$: the probability that two independent root queries both return
true is asymptotic to $q(1)^2\sim(2D)^{-1}$. A collision estimate that merely
tends to zero would not control that second moment.

Throughout this section, $G$ is any graph satisfying
Proposition~\ref{prop:prefix}, and all constants and error estimates are
uniform over such $G$. We work in the independent unfolding. The proof has
two parts. First, the extra edge between the marked vertices makes their
path union rare as a graph embedding. Second, reaching a long prescribed
path requires many earlier candidate queries to fail. We retain the order
of the selected priorities when combining these two bounds.

\subsection{Patterns and their embeddings}

The paths in this section are paths in the \emph{call forest}, not paths
in the formal graph $\Gamma$. Consider the union of the ancestral paths
to the birth calls of two marked formal vertices $a,b$. When the paths
belong to the same rooted tree, let $h$ be their common trunk length and
let $\ell_1,\ell_2$ be the two lengths after the trunk. Length counts
non-root calls. When the paths belong to distinct trees, set $h=0$ and
let $\ell_1,\ell_2$ be their individual lengths. In both cases put
\[
 \ell=h+\ell_1+\ell_2,
 \qquad r=\text{the number of roots used}\in\{1,2\}.
\]
Zero lengths are allowed. A mark born at a root specifies one of its two
formal endpoints.

A \emph{pattern} records these rooted paths, their marks, and the choice
of focus edge at each child step. The named focus edges from
Section~\ref{sec:coupling} give at most two choices at a step. The choices
of roots, the order of the two marked paths, and any marked root endpoints
contribute only an absolute constant. Consequently, for each triple
$(h,\ell_1,\ell_2)$ there are at most
\begin{equation}\label{eq:pattern-count}
 C\,2^\ell
\end{equation}
patterns, for an absolute constant $C$. At a fork the two children each
choose a focus edge; their third vertices will be determined by an
embedding. Thus there is no further length-dependent multiplicity.
Patterns that cannot occur under the query rule may be discarded.



The formal graph $J$ of a pattern starts with $r$ disjoint root edges and
adds two edges and one vertex at each non-root call. In particular,
\[
 v(J)=2r+\ell,\qquad e(J)=r+2\ell.
\]
We consider only patterns whose marks $a,b$ are distinct and nonadjacent
in $J$, and write $J^+=J+ab$. An embedding below is an injective map of
the formal vertices into $V(G)$ that maps every edge of $J^+$ to an edge
of $G$ and respects the named endpoints of each root.

\begin{lemma}[Embedding an extra edge]\label{lem:embeddings}
For every such pattern with $r\in\{1,2\}$ roots and $\ell$ non-root calls,
the number of embeddings of $J^+$ into $G$ is at most
\begin{equation}\label{eq:embedding-bound}
 (2m)^r(2D)^\ell D^{-B}
\end{equation}
for all sufficiently large $n$, uniformly over the pattern and $G$.
\end{lemma}

\begin{proof}
Write $\kappa_n=(1+\log n)^{C_0}$ for the common factor in the two template
bounds of Proposition~\ref{prop:prefix}. Recall that $L=100$, $B=50$,
$p\sim n^{-1/2+\epsilon}$, and $D=np^2\sim n^{2\epsilon}$, where
$\epsilon=1/2000$.

\paragraph{Fewer than $L$ calls.}
Suppose $\ell<L$. Regard $J^+$ as a template with no distinguished
vertices. Order all root vertices first, ordering the endpoints of each
root edge, and then add the other vertices in any parent-first order of
their birth calls. Before adding $ab$, each vertex has at most two
neighbors preceding it. Moreover, the first vertex of any nonempty
vertex subset has no preceding neighbor within that subset. An induced
subgraph on $v>0$ vertices therefore has at most $2v-2$ edges from $J$,
and at most $2v-1\leq 2v$ from $J^+$. Every subgraph satisfies the same
weaker bound $e\leq 2v$. Since $p<1$ and $D>1$ for large $n$, its scaling
is at least
\[
 n^v p^{2v}=D^v\geq1.
\]
The empty subgraph has scaling one. The first template bound gives
\[
 \#\{J^+\hookrightarrow G\}
 \leq \kappa_n n^{2r+\ell}p^{r+2\ell+1}
 =\kappa_n(2m)^rD^\ell p.
\]
Here the template has at most $4+(L-1)=103\leq H_0$ vertices. The
inequality
\[
 pD^B=n^{-1/2+(2B+1)\epsilon+o(1)}
      =n^{-899/2000+o(1)}
\]
shows that $\kappa_np\leq D^{-B}$ for large $n$. This proves
\eqref{eq:embedding-bound} in the short case.

\paragraph{The last $L$ births.}
Suppose $\ell\geq L$. Choose any parent-first ordering of the non-root
calls, with all roots placed first, and let $Z$ be the set of the last
$L$ vertices born. The labels of the preceding vertices can be assigned
in at most
\begin{equation}\label{eq:preceding-labels}
 (2m)^r(2D)^{\ell-L}
\end{equation}
ways. Indeed, there are $2m$ choices for each oriented root, and at most
$2D$ common neighbors of the two endpoints of each subsequent attachment
edge. Discarding restrictions against repeated labels only increases this
bound. We now fix an injective assignment of these preceding labels.

The birth call of every vertex in $Z$ is an ancestor of the birth call
of a marked vertex in $Z$. To see this, every call in the pattern lies
on a path to a marked birth call. Once that path enters the final suffix,
all its subsequent calls lie in the suffix, by the parent-first order.
In particular, at least one endpoint of $ab$ belongs to $Z$.
Figure~\ref{fig:call-paths} illustrates how such a suffix can meet both arms.

\begin{figure}[htbp]
\centering
\begin{tikzpicture}[x=1cm,y=1cm,
 call/.style={circle,draw=black!65,fill=white,minimum size=6mm,inner sep=0pt},
 last/.style={call,fill=linkblue!18,draw=linkblue},
 every label/.style={font=\small}]
\node[call,label=above:root] (o) at (0,0) {$0$};
\node[call] (t1) at (0,-.7) {$1$};
\node[call] (t2) at (0,-1.4) {$2$};
\node[call] (a1) at (-1.4,-2.1) {$3$};
\node[last,label=below:$\operatorname{birth}(a)$] (a) at (-1.4,-3.5) {$5$};
\node[last] (b1) at (1.4,-2.1) {$4$};
\node[last] (b2) at (1.4,-2.8) {$6$};
\node[last,label=below:$\operatorname{birth}(b)$] (b) at (1.4,-3.5) {$7$};
\draw[black!65] (o)--(t1)--(t2)--(a1)--(a);
\draw[black!65] (t2)--(b1)--(b2)--(b);
\node at (-.85,-.7) {$h=2$};
\node at (-2.45,-2.65) {$\ell_1=2$};
\node at (2.45,-2.65) {$\ell_2=3$};
\end{tikzpicture}
\caption{A union of two ancestral call paths. Nodes are calls and lines
are parent--child relations, not edges of $G$ or $\Gamma$.
Each non-root call introduces one formal vertex.
The numbers give one parent-first ordering; the last four non-root calls
are shaded. Such a final suffix can occupy both arms and need not be
connected. This small example illustrates the final block used below,
where the proof takes $L=100$.}
\label{fig:call-paths}
\end{figure}

Form a template $(H,I)$ as follows. Its free vertices are $Z$. Retain the
two backward edges of each vertex of $Z$, meaning its edges to the
endpoints of its attachment edge, and retain $ab$. The distinguished set
$I$ consists of all preceding vertices needed as endpoints of these
edges. Include no edges with both endpoints in $I$. We may discard such
constraints because we seek an upper bound on extensions. Thus $I$ is
independent in the template, whether or not its prescribed images have
edges between them in $G$.

The backward edges use at most $2L$ boundary vertices, and the extra edge
uses at most one further boundary vertex. Hence
\[
 |I|\leq2L+1,\qquad v(H)\leq3L+1\leq H_0.
\]
The $2L$ backward edges are distinct, and $ab$ is not one of them.
Every retained edge touches $Z$, so
\begin{equation}\label{eq:terminal-scaling}
 e(H)=2L+1,\qquad
 S(H,I)=n^Lp^{2L+1}=D^Lp
       =n^{-799/2000+o(1)}<1.
\end{equation}

We verify every completion condition needed for the second template
bound. If $U\subsetneq Z$, then
\begin{equation}\label{eq:proper-free-edges}
 e\bigl(H[I\cup U]\bigr)\leq2|U|.
\end{equation}
Indeed, assign each backward edge to its newer endpoint in $Z$. In
$H[I\cup U]$ there are at most two such edges per vertex of $U$.
If \eqref{eq:proper-free-edges} failed, the induced graph would retain
$ab$ and both backward edges of every vertex of $U$. Retaining $ab$
puts every free mark in $U$. Retaining both backward edges of a vertex
also forces the vertex born at its immediate parent call into $U$,
whenever that vertex belongs to $Z$. Indeed, both focus edges at a
non-root parent contain its newly born vertex, so the child's attachment
edge contains that vertex. Repeating this argument backward along each
marked path includes every vertex of $Z$, by the ancestor property just
proved. This would give $U=Z$, a contradiction.

For a proper $U$, \eqref{eq:proper-free-edges} implies
\[
 S\bigl(H[I\cup U],I\bigr)
 =n^{|U|}p^{e(H[I\cup U])}\geq D^{|U|}\geq1.
\]
Every intermediate set of distinguished vertices has the form
$J_0=I\cup U$, with $U\subseteq Z$. Its completion scaling is
\[
 n^{v(H)-|J_0|}p^{e(H)-e(H[J_0])}
 =\frac{S(H,I)}{S(H[J_0],I)}.
\]
For proper $U$ the numerator is less than one by
\eqref{eq:terminal-scaling}, while the denominator is at least one.
For $U=Z$ the ratio equals one. All the required completion scalings are
therefore at most one.

The second template bound now gives at most $\kappa_n$ extensions for the
labels of $Z$. In applying it we may ignore avoidance of previously
assigned vertices outside $I$, again only enlarging the count. Since
$L-B=50>0$ and $D\sim n^{1/1000}$,
\[
 \kappa_n\leq(2D)^L D^{-B}
\]
for all sufficiently large $n$. Multiplying this by
\eqref{eq:preceding-labels} proves \eqref{eq:embedding-bound}.
The threshold for $n$ is uniform: the template order and its
polylogarithmic exponent are fixed independently of the path lengths.
\end{proof}

The embedding bound supplies a factor $D^{-B}$ even when the paths are
arbitrarily long. It remains to sum over their lengths. Counting paths
alone would lose this factor, since a call can have order $D$ children.
The next estimate uses the failures required to reach a prescribed
child, together with the decreasing order of priorities.

\subsection{The probability of visiting prescribed paths}

Fix a pattern and an embedding of $J^+$ \emph{before} exposing any
priorities. The probability that its $r$ used roots have the prescribed
ordered labels is $(2m)^{-r}$. Conditional on this root event, each step
of the pattern specifies one candidate occurrence in a fixed potential
call tree: its focus edge and third vertex determine its triangle.
If these specifications cannot define distinct required children under
the query rule, their visitation probability is zero and they can be
discarded. Otherwise, write $u_y$ for the priority of the selected child
occurrence corresponding to each of the $\ell$ non-root calls $y$.
These are distinct independent uniform coordinates of the unfolding.

Set
\begin{equation}\label{eq:visitation-weights}
 b(t)=q(t)^2=\frac{1}{1+2Dt},\qquad
 C_D=1+2D,\qquad
 \mu=\int_0^1b(t)\,dt=\frac{\log(1+2D)}{2D}.
\end{equation}

\begin{lemma}[Visiting an embedded pattern]\label{lem:visitation}
Conditional on the prescribed ordered root labels, the probability that
all calls in the embedded pattern are visited is at most
\begin{equation}\label{eq:visitation-bound}
 C_D^3\frac{(2\mu)^\ell}{h!\,\ell_1!\,\ell_2!}.
\end{equation}
This includes zero-length segments, with the convention $0!=1$.
\end{lemma}

\begin{proof}
Condition only on the selected priority coordinates $(u_y)$, not on
actual visitation, return values, or previously exposed lists. For all
required calls to be visited, these priorities must decrease along every
required path.

Consider a required non-root call $x$ that has one or two required
children. Its focus is a pair $\{f,g\}$, and its omitted parent triangle
is $T$. Let $t_x$ be the largest priority of its required children.
Define $O_x$ prospectively to be the event that every off-pattern
candidate at $x$ with priority below $t_x$ has a child query returning
false, when that subtree is evaluated hypothetically. Reaching all the
required children necessitates $O_x$. In particular,
\begin{equation}\label{eq:necessary-offpattern}
 \{\text{all required calls are visited}\}
 \subseteq
 \{\text{priorities decrease along the required paths}\}
 \cap\bigcap_x O_x.
\end{equation}
If a required call has two required children, reaching the later child
also requires the earlier selected child to return false. We discard
this additional condition in \eqref{eq:necessary-offpattern}.

Each $O_x$ depends only on subtrees rooted at off-pattern children of
$x$. These subtrees contain none of the selected priority coordinates
and are disjoint for different required calls. In particular, when one
required call is an ancestor of another, the path between them starts
through a selected child, so it does not enter an off-pattern subtree.
The events $O_x$ are therefore independent conditional on $(u_y)$.
This assertion concerns the prospective independent unfolding.

For an off-pattern candidate $S$ attached through focus edge $e$, the
probability that it is not a successful child before $t_x$ is exactly
$1-I_{e,S}(t_x)$. By the product identities of
Proposition~\ref{prop:products},
\begin{equation}\label{eq:offpattern-product}
 \Prb(O_x\mid (u_y))
 =\frac{q_{f,T}(t_x)q_{g,T}(t_x)}
 {\displaystyle\prod_{S\text{ selected at }x}
       (1-I_{e(S),S}(t_x))},
\end{equation}
where $e(S)$ is the unique focus edge contained in $S$.
The denominator removes the one or two selected candidate factors from
the complete product. The product estimate~\eqref{eq:off-pattern-bound}
therefore gives, uniformly in the types and in $t_x\in[0,1]$,
\begin{equation}\label{eq:offpattern-bound}
 \Prb(O_x\mid (u_y))\leq 2b(t_x)
\end{equation}
for all sufficiently large $n$.

We next distribute these factors among the selected child priorities.
At a call with one required child of priority $u$, the bound is $2b(u)$.
At a call with two required children of priorities $u,v$, we use
\[
 \frac{b(\max\{u,v\})}{b(u)b(v)}
 =\frac{1}{b(\min\{u,v\})}\leq C_D.
\]
There is at most one such non-root fork. We ignore the off-pattern
requirements at roots entirely. There are at most two required child
steps leaving roots in total, and inserting their missing factors
$b(u)\geq C_D^{-1}$ costs at most $C_D^2$.
The number of non-root calls with required children is at most $\ell$.
Combining independence, \eqref{eq:necessary-offpattern}, and
\eqref{eq:offpattern-bound}, we obtain the conditional bound
\begin{equation}\label{eq:conditional-visitation}
 C_D^3 2^\ell\prod_y b(u_y)
\end{equation}
when the priorities decrease along the required paths, and zero
otherwise. Every selected priority now carries one factor $b$.
This also covers root marks and empty arms, which introduce no new
priority coordinates.

The selected priorities have joint Lebesgue density one. In integrating
\eqref{eq:conditional-visitation}, retain the order constraints within
the trunk and within each arm, but discard constraints joining different
segments. For a segment of length $j$, symmetry of the product integrand
gives
\[
 \int_{1>u_1>\cdots>u_j>0}\prod_{i=1}^j b(u_i)\,du_1\cdots du_j
 =\frac{\mu^j}{j!}.
\]
The three segments use disjoint coordinates. Their contributions
multiply, proving \eqref{eq:visitation-bound}. Priority ties have
probability zero.
\end{proof}

\subsection{Summing all witnesses}

\begin{proposition}[Uniform collision bound]\label{prop:collision}
Fix $k\in\{1,2\}$ and sample $k$ independent uniform oriented edges of
$G$, with replacement. Execute the independent root queries at threshold
$1$, executing every root regardless of earlier outputs. Let
$\mathcal C_k$ be the event that an actual triangle type is repeated among
all exposed candidate lists, including candidates not traversed. Then
\begin{equation}\label{eq:collision-bound}
 \Prb_{\mathrm{ind}}(\mathcal C_k)
 =O(D^{15-B}+n^{-1})
 =O(D^{-35}+n^{-1})
 =o(D^{-1}).
\end{equation}
The bound is uniform over the graphs in Proposition~\ref{prop:prefix}.
Consequently, by Lemma~\ref{lem:coupling}, the probabilities that all root
queries return true in the finite and independent experiments differ
by at most the same bound.
\end{proposition}

\begin{proof}
Lemma~\ref{lem:witness} reduces a collision to initial root overlap,
which has probability $O(n^{-1})$, or to an injectively labeled marked
path union with an extra edge. For each fixed length triple, multiply
the pattern count \eqref{eq:pattern-count}, the embedding bound
\eqref{eq:embedding-bound}, the root-matching probability $(2m)^{-r}$,
and the visitation bound \eqref{eq:visitation-bound}. The result is at
most an absolute constant times
\[
 D^{-B}C_D^3
 \frac{(8D\mu)^{h+\ell_1+\ell_2}}{h!\,\ell_1!\,\ell_2!}.
\]
For two-tree patterns $h=0$; allowing all triples of nonnegative
integers only enlarges the sum. A union bound over patterns and
embeddings therefore bounds the witness probability by
\begin{align*}
 O\left(D^{-B}C_D^3
   \sum_{h,\ell_1,\ell_2\geq0}
   \frac{(8D\mu)^{h+\ell_1+\ell_2}}{h!\,\ell_1!\,\ell_2!}\right)
 &=O\left(D^{-B}C_D^3\exp(24D\mu)\right)\\
 &=O\left(D^{-B}C_D^{15}\right)
  =O(D^{15-B}).
\end{align*}
All terms are nonnegative; the countable union bound applies to the
almost surely finite evaluations established in
Section~\ref{sec:queries}. The identity
$24D\mu=12\log C_D$ accounts for the exponent $15$.
Finally, $B=50$ and $D\sim n^{1/1000}$ imply
$D^{-35}+n^{-1}=o(D^{-1})$, proving \eqref{eq:collision-bound}.
\end{proof}

We have thus compared the finite continuation with independent root
queries at an error smaller than their two-root success probability.
The first and second terminal moments can now be read directly from
these root-query probabilities.
